\section{AIME I 2003}
        \begin{enumerate}
        	\item (\textit{AIME I 2003}) Given that


 $\frac{((3!)!)!}{3!} = k \cdot n!,$

 where $k$ and $n$ are positive integers and $n$ is as large as possible, find $k + n.$


 
 \textbf{Answer: }839



\textbf{Solution}

Note that
 \[{{\left((3!)!\right)!}\over{3!}}= {{(6!)!}\over{6}}={{720!}\over6}={{720\cdot719!}\over6}=120\cdot719!.\]Because $120\cdot719!<720!$, we can conclude that $n < 720$. Thus, the maximum value of $n$ is $719$. The requested value of $k+n$ is therefore $120+719=\boxed{839}$.

 


	\item (\textit{AIME I 2003}) One hundred concentric circles with radii $1, 2, 3, \dots, 100$ are drawn in a plane. The interior of the circle of radius 1 is colored red, and each region bounded by consecutive circles is colored either red or green, with no two adjacent regions the same color. The ratio of the total area of the green regions to the area of the circle of radius 100 can be expressed as $m/n,$ where $m$ and $n$ are relatively prime positive integers. Find $m + n.$


 
 \textbf{Answer: }301



\textbf{Solution 1}

To get the green area, we can color all the circles of radius $100$ or below green, then color all those with radius $99$ or below red, then color all those with radius $98$ or below green, and so forth.  This amounts to adding the area of the circle of radius $100$, but subtracting the circle of radius $99$, then adding the circle of radius $98$, and so forth.

 The total green area is thus given by $100^{2} \pi - 99^{2} \pi + 98^{2} \pi - \ldots - 1^{2} \pi$, while the total area is given by $100^{2} \pi$, so the ratio is
 \[\frac{100^{2}\pi - 99^{2}\pi + 98^{2}\pi - \ldots - 1^{2}\pi}{100^{2}\pi}\]

 For any $a$, $a^{2}-(a-1)^{2}=a^{2}-(a^{2}-2a+1)=2a-1$.  We can cancel the \href{https://artofproblemsolving.com/wiki/index.php?title=Divisor}{ factor} of \href{https://artofproblemsolving.com/wiki/index.php?title=Pi}{pi} from the \href{https://artofproblemsolving.com/wiki/index.php?title=Numerator}{numerator} and \href{https://artofproblemsolving.com/wiki/index.php?title=Denominator}{denominator} and simplify the ratio to 

 
 \[\frac{(2\cdot100 - 1)+(2\cdot98 - 1) + \ldots + (2\cdot 2 - 1)}{100^{2}} = \frac{2\cdot(100 + 98 + \ldots + 2) - 50}{100^2}.\] 

 Using the formula for the sum of an \href{https://artofproblemsolving.com/wiki/index.php?title=Arithmetic\_series}{arithmetic series}, we see that this is equal to 

 
 \[\frac{2(50)(51)-50}{100^{2}}=\frac{50(101)}{100^{2}}=\frac{101}{200},\]

 so the answer is $101 + 200 =\boxed{301}$.

 Alternatively, we can determine a pattern through trial-and-error using smaller numbers. 

 
                \begin{itemize}
                    \item For $2$ circles, the ratio is $3/4$.\item For $4$ circles, the ratio is $5/8$.\item For $6$ circles, the ratio is $7/12$.\item For $8$ circles, the ratio is $9/16$.
                \end{itemize}


                Now the pattern for each ratio is clear.  Given $x$ circles, the ratio is $\frac{x+1}{2x}$.For the $100$ circle case (which is what this problem is), $x=100$, and the ratio is $\frac{101}{200}$.

 


 Also, using the difference of squares, the expression simplifies to $\frac{100 + 99 + 98 + 97 + ... + 1}{100^2}$. We can easily determine the sum with $\frac{100(101)}{2} = 5050$. Simplifying gives us $\frac{5050}{100^2} = \frac{101}{200}$ and the answer is $101 + 200 =\boxed{301}$.

 



\textbf{Solution 2 (synergy)}

We want to find $\frac{\sum\limits_{n=1}^{50} (4n-1)\pi}{10000\pi}=\frac{\sum\limits_{n=1}^{50} (4n-1)}{10000}=\frac{(\sum\limits_{n=1}^{50} (4n) )-50}{10000}=\frac{101}{200} \rightarrow 101+200=\boxed{301}$

 



\textbf{Solution 3 (Alcumus)}

The sum of the areas of the green regions is

 
 \[\left[(2^2-1^2)+(4^2-3^2)+(6^2-5^2)+\cdots+(100^2-99^2)\right]\pi\]
 \[=\left[(2+1)+(4+3)+(6+5)+\cdots+(100+99)\right]\pi\]
 \[={1\over2}\cdot100\cdot101\pi.\]

 Thus the desired ratio is
 \[{1\over2}\cdot{{100\cdot101\pi}\over{100^2\pi}}={101\over200},\]and $m+n=\boxed{301}$.

 


	\item (\textit{AIME I 2003}) Let the set $\mathcal{S} = \{8, 5, 1, 13, 34, 3, 21, 2\}.$ Susan makes a list as follows: for each two-element subset of $\mathcal{S},$ she writes on her list the greater of the set's two elements. Find the sum of the numbers on the list.


 
 \textbf{Answer: }484



\textbf{Solution}

Order the numbers in the set from greatest to least to reduce error: $\{34, 21, 13, 8, 5, 3, 2, 1\}.$Each \href{https://artofproblemsolving.com/wiki/index.php?title=Element}{element} of the \href{https://artofproblemsolving.com/wiki/index.php?title=Set}{set} will appear in $7$ two-element \href{https://artofproblemsolving.com/wiki/index.php?title=Subset}{subsets}, once with each other number.

 
                \begin{itemize}
                    \item $34$ will be the greater number in $7$ subsets.\item $21$ will be the greater number in $6$ subsets.\item $13$ will be the greater number in $5$ subsets.\item $8$ will be the greater number in $4$ subsets.\item $5$ will be the greater number in $3$ subsets.\item $3$ will be the greater number in $2$ subsets.\item $2$ will be the greater number in $1$ subsets.\item $1$ will be the greater number in $0$ subsets.
                \end{itemize}


                Therefore the desired sum is $34\cdot7+21\cdot6+13\cdot5+8\cdot4+5\cdot3+3 \cdot2+2\cdot1+1\cdot0=\boxed{484}$.

 Note: Note that $7+6+5+4+3+2+1=\binom{8}{2}$, so we have counted all the possible cases.

 


 Interesting note: All the numbers in the set are Fibonacci numbers

 



\textbf{Solution}

Thinking of this problem algorithmically, one can "sort" the array to give:
 \[{1, 2, 3, 5, 8, 13, 21, 34}\]

 Now, notice that when we consider different pairs, we are only going to fixate one element and look at the all of the next elements in the array, basically the whole $j = i + 1$ shebang. Then, we see that if we set the sum of the whole array to $x,$ we get out answer to be

 
 \[(x-1) + (x-3) + (x-6) + (x-11) + (x-19) + (x-32) + (x-53) = 7x - 125\]

 Finding $7x$ isn't hard, and we see that it is equal to $609$:

 
 \[609 - 125 = \boxed{484}\]

 


	\item (\textit{AIME I 2003}) Given that $\log_{10} \sin x + \log_{10} \cos x = -1$ and that $\log_{10} (\sin x + \cos x) = \frac{1}{2} (\log_{10} n - 1),$ find $n.$


 
 \textbf{Answer: }012



\textbf{Solution 1}

Using the properties of \href{https://artofproblemsolving.com/wiki/index.php?title=Logarithm}{logarithms}, we can simplify the first equation to $\log_{10} \sin x + \log_{10} \cos x = \log_{10}(\sin x \cos x) = -1$. Therefore, 
 \[\sin x \cos x = \frac{1}{10}.\qquad (*)\]

 Now, manipulate the second equation.
 \begin{align*} \log_{10} (\sin x + \cos x) &= \frac{1}{2}(\log_{10} n - \log_{10} 10) \\ \log_{10} (\sin x + \cos x) &= \left(\log_{10} \sqrt{\frac{n}{10}}\right) \\ \sin x + \cos x &= \sqrt{\frac{n}{10}} \\ (\sin x + \cos x)^{2} &= \left(\sqrt{\frac{n}{10}}\right)^2 \\ \sin^2 x + \cos^2 x +2 \sin x \cos x &= \frac{n}{10} \\ \end{align*}


 By the Pythagorean identities, $\sin ^2 x + \cos ^2 x = 1$, and we can substitute the value for $\sin x \cos x$ from $(*)$. $1 + 2\left(\frac{1}{10}\right) = \frac{n}{10} \Longrightarrow n = \boxed{012}$.

 



\textbf{Solution 2}

Examining the first equation, we simplify as the following:
 \[\log_{10} \sin x \cos x = -1\]
 \[\implies \sin x \cos x = \frac{1}{10}\]

 With this in mind, examining the second equation, we may simplify as the following (utilizing logarithm properties):
 \[\log_{10} (\sin x + \cos x) = \frac{1}{2} (\log_{10} n - \log_{10} 10)\]
 \[\implies \log_{10} (\sin x + \cos x) = \frac{1}{2} (\log_{10} \frac{n}{10})\]
 \[\implies \log_{10} (\sin x + \cos x) = \log_{10} \sqrt{\frac{n}{10}}\]

 From here, we may divide both sides by $\log_{10} (\sin x + \cos x)$ and then proceed with the change-of-base logarithm property:
 \[1 = \frac{\log_{10} \sqrt{\frac{n}{10}}}{\log_{10} (\sin x + \cos x)}\]
 \[\implies 1 = \log_{\sin x + \cos x} \sqrt{\frac{n}{10}}\]

 Thus, exponentiating both sides results in $\sin x + \cos x = \sqrt{\frac{n}{10}}$. Squaring both sides gives us 
 \[\sin^2 x + 2\sin x \cos x + \cos^2 x = \frac{n}{10}\]

 Via the Pythagorean Identity, $\sin^2 x + \cos^2 x = 1$ and $2\sin x \cos x$ is simply $\frac{1}{5}$, via substitution. Thus, substituting these results into the current equation:
 \[1 + \frac{1}{5} = \frac{n}{10}\]
 \[\implies \frac{6}{5} = \frac{n}{10}\]

 Using simple cross-multiplication techniques, we have $5n = 60$, and thus $\boxed{n = 012}$.~ nikenissan

 


 



\textbf{Solution 3}

By the first equation, we get that $\sin(x)*\cos(x)=10^{-1}$. We can let $\sin(x)=a$, $\cos(x)=b$. Thus $ab=\frac{1}{10}$. By the identity $\sin^2x+\cos^2x=1$, we get that $a^2+b^2=1$. Solving this, we get $a+b=\sqrt{\frac{12}{10}}$. So we have

 
 \[\log\left(\sqrt{\frac{12}{10}}\right)=\frac12(\log(n)-1)\]
 \[2\log\left(\sqrt{\frac{12}{10}}\right)=\log(n)-1\]
 \[\log\left(\frac{12}{10}\right)+1=\log(n)\]
 \[\log\left(\frac{12}{10}\right)+\log(10)=\log(n)\]
 \[\log\left(\frac{12}{10}\times 10\right)=\log(12)=\log(n)\]

 From here it is obvious that $\boxed{n=012}$.

 



\textbf{Solution 4}

Let $\log{x} = \log_{10}{x}.$ Through basic log properties, we see that $\log{a} + \log{b} = \log{(ab)}.$ Thus, we see that $\log{(\sin{x})} + \log{(\cos{x})} = \log{(\sin{x}\cos{x})} = -1.$ Simplifying, we get:

 \begin{align*}\log{(\sin{x}\cos{x})} &= -1 \\\sin{x}\cos{x} &= 10^{-1} = \frac{1}{10}\end{align*}


 Next, we can manipulate the second equation to get:

 \begin{align*}\log{(\sin{x} + \cos{x})} &= \frac{1}{2}(\log{n}-1) \\2\log{(\sin{x} + \cos{x})} &= \log{n}-1 \\\log{(\sin{x} + \cos{x})^2} + 1 &= \log{n}\end{align*}


 Expanding $(\sin{x} + \cos{x})^2,$ we get:

 \begin{align*}\log{(\sin^2{x} + \cos^2{x} + 2\sin{x}\cos{x})} + 1 &= \log{n} \\\log{(1 + 2\sin{x}\cos{x})} + 1 &= \log{n} \\\log{(1 + \frac{2}{10})} + \log{10} &= \log{n} \\\log{(\frac{12}{10} \cdot 10)} = \log{n} \\\log{12} = \log{n}\end{align*}


 Finally, we see that $n = \boxed{012}.$

 


	\item (\textit{AIME I 2003}) Consider the set of points that are inside or within one unit of a rectangular parallelepiped (box) that measures 3 by 4 by 5 units. Given that the volume of this set is $\frac{m + n \pi}{p},$ where $m, n,$ and $p$ are positive integers, and $n$ and $p$ are relatively prime, find $m + n + p.$


 
 \textbf{Answer: }505



\textbf{Solution}


                    \begin{figure}[H]
                        \centering
                        \includegraphics[width=0.5\textwidth]{Solution/AIME_I_2003_5_Solution_Figure_1.png}
                    \end{figure}
                    

 The set can be broken into several parts: the big $3\times 4 \times 5$ parallelepiped, $6$ external parallelepipeds that each share a face with the large parallelepiped and have a height of $1$, the $1/8$ \href{https://artofproblemsolving.com/wiki/index.php?title=Sphere}{spheres} (one centered at each \href{https://artofproblemsolving.com/wiki/index.php?title=Vertex}{vertex} of the large parallelepiped), and the $1/4$ \href{https://artofproblemsolving.com/wiki/index.php?title=Cylinder}{cylinders} connecting each adjacent pair of spheres.

 
                \begin{itemize}
                    \item The volume of the parallelepiped is $3 \times 4 \times 5 = 60$ cubic units.\item The volume of the external parallelepipeds is $2(3 \times 4 \times 1)+2(3 \times 5 \times 1 )+2(4 \times 5 \times 1)=94$.\item There are $8$ of the $1/8$ spheres, each of radius $1$. Together, their volume is $\frac{4}{3}\pi$.\item There are $12$ of the $1/4$ cylinders, so $3$ complete cylinders can be formed. Their volumes are $3\pi$, $4\pi$, and $5\pi$, adding up to $12\pi$.
                \end{itemize}


                The combined volume of these parts is $60+94+\frac{4}{3}\pi+12\pi = \frac{462+40\pi}{3}$. Thus, the answer is $m+n+p = 462+40+3 = \boxed{505}$.

 


	\item (\textit{AIME I 2003}) The sum of the areas of all triangles whose vertices are also vertices of a 1 by 1 by 1 cube is $m + \sqrt{n} + \sqrt{p},$ where $m, n,$ and $p$ are integers. Find $m + n + p.$


 
 \textbf{Answer: }348



\textbf{Solution}


                    \begin{figure}[H]
                        \centering
                        \includegraphics[width=0.5\textwidth]{Solution/AIME_I_2003_6_Solution_Figure_1.png}
                    \end{figure}
                    

 
                    \begin{figure}[H]
                        \centering
                        \includegraphics[width=0.5\textwidth]{Solution/AIME_I_2003_6_Solution_Figure_2.png}
                    \end{figure}
                    

 
                    \begin{figure}[H]
                        \centering
                        \includegraphics[width=0.5\textwidth]{Solution/AIME_I_2003_6_Solution_Figure_3.png}
                    \end{figure}
                    

 

 Since there are $8$ \href{https://artofproblemsolving.com/wiki/index.php?title=Vertex}{ vertices} of a \href{https://artofproblemsolving.com/wiki/index.php?title=Cube\_(geometry)}{ cube}, there are ${8 \choose 3} = 56$ total \href{https://artofproblemsolving.com/wiki/index.php?title=Triangle}{triangles} to consider.  They fall into three categories: there are those which are entirely contained within a single \href{https://artofproblemsolving.com/wiki/index.php?title=Face}{face} of the cube (whose sides are two \href{https://artofproblemsolving.com/wiki/index.php?title=Edge}{edges} and one face \href{https://artofproblemsolving.com/wiki/index.php?title=Diagonal}{diagonal}), those which lie in a \href{https://artofproblemsolving.com/wiki/index.php?title=Plane}{plane} \href{https://artofproblemsolving.com/wiki/index.php?title=Perpendicular}{perpendicular} to one face of the cube (whose sides are one edge, one face diagonal and one space diagonal of the cube) and those which lie in a plane \href{https://artofproblemsolving.com/wiki/index.php?title=Oblique&action=edit&redlink=1}{oblique} to the edges of the cube, whose sides are three face diagonals of the cube.  

 Each face of the cube contains ${4\choose 3} = 4$ triangles of the first type, and there are $6$ faces, so there are $24$ triangles of the first type. Each of these is a \href{https://artofproblemsolving.com/wiki/index.php?title=Right\_triangle}{right triangle} with legs of length $1$, so each triangle of the first type has area $\frac 12$.

 Each edge of the cube is a side of exactly $2$ of the triangles of the second type, and there are $12$ edges, so there are $24$ triangles of the second type. Each of these is a right triangle with legs of length $1$ and $\sqrt 2$, so each triangle of the second type has area $\frac{\sqrt{2}}{2}$.

 Each vertex of the cube is associated with exactly one triangle of the third type (whose vertices are its three neighbors), and there are $8$ vertices of the cube, so there are $8$ triangles of the third type. Each of the these is an \href{https://artofproblemsolving.com/wiki/index.php?title=Equilateral\_triangle}{equilateral triangle} with sides of length $\sqrt 2$, so each triangle of the third type has area $\frac{\sqrt 3}2$.

 Thus the total area of all these triangles is $24 \cdot \frac12 + 24\cdot\frac{\sqrt2}2 + 8\cdot\frac{\sqrt3}2 = 12 + 12\sqrt2 + 4\sqrt3 = 12 + \sqrt{288} + \sqrt{48}$ and the answer is $12 + 288 + 48 = \boxed{348}$.

 


	\item (\textit{AIME I 2003}) Point $B$ is on $\overline{AC}$ with $AB = 9$ and $BC = 21.$ Point $D$ is not on $\overline{AC}$ so that $AD = CD,$ and $AD$ and $BD$ are integers. Let $s$ be the sum of all possible perimeters of $\triangle ACD.$ Find $s.$


 
 \textbf{Answer: }380



\textbf{Solution 1 (Pythagorean Theorem)}


                    \begin{figure}[H]
                        \centering
                        \includegraphics[width=0.5\textwidth]{Solution/AIME_I_2003_7_Solution_Figure_1.png}
                    \end{figure}
                    

 Denote the height of $\triangle ACD$ as $h$, $x = AD = CD$, and $y = BD$. Using the \href{https://artofproblemsolving.com/wiki/index.php?title=Pythagorean\_theorem}{Pythagorean theorem}, we find that $h^2 = y^2 - 6^2$ and $h^2 = x^2 - 15^2$. Thus, $y^2 - 36 = x^2 - 225 \Longrightarrow x^2 - y^2 = 189$. The LHS is \href{https://artofproblemsolving.com/wiki/index.php?title=Difference\_of\_squares}{difference of squares}, so $(x + y)(x - y) = 189$. As both $x,\ y$ are integers, $x+y,\ x-y$ must be integral divisors of $189$. 

 The pairs of divisors of $189$ are $(1,189)\ (3,63)\ (7,27)\ (9,21)$. This yields the four potential sets for $(x,y)$ as $(95,94)\ (33,30)\ (17,10)\ (15,6)$. The last is not a possibility since it simply \href{https://artofproblemsolving.com/wiki/index.php?title=Degenerate}{degenerates} into a \href{https://artofproblemsolving.com/wiki/index.php?title=Line}{line}. The sum of the three possible perimeters of $\triangle ACD$ is equal to $3(AC) + 2(x_1 + x_2 + x_3) = 90 + 2(95 + 33 + 17) = \boxed{380}$.

 \textbf{Note} One can also conclude $(15,6)$ doesn't work as it violates the triangle inequality.

 



\textbf{Solution 2 (Stewart's Theorem)}

Let $AD=c$ and $BD=d$, then by \href{https://artofproblemsolving.com/wiki/index.php?title=Stewart\%27s\_Theorem}{Stewart's Theorem} we have:

 $30d^2+21*9*30=9c^2+21c^2=30c^2$. After simplifying:

 $d^2-c^2=189$.

 The solution follows as above.

 



\textbf{Solution 3 (Law of Cosines)}

Drop an altitude from point $D$ to side $AC$. Let the intersection point be $E$. Since triangle $ADC$ is isosceles, AE is half of $AC$, or $15$. Then, label side AD as $x$. Since $AED$ is a right triangle, you can figure out $\cos A$ with adjacent divided by hypotenuse, which in this case is $AE$ divided by $x$, or $\frac{15}{x}$. Now we apply law of cosines. Label $BD$ as $y$. Applying law  of cosines, $y^2 = x^2+9^2- 2 \cdot x \cdot 9 \cdot \cos A$. Since $\cos A$ is equal to $\frac{15}{x}$, $y^2 = x^2+9^2- 2 \cdot x \cdot 9 \cdot \frac{15}{x}$, which can be simplified to $x^2-y^2=189$. The solution proceeds as the first solution does.

 -intelligence\_20

 


	\item (\textit{AIME I 2003}) In an increasing sequence of four positive integers, the first three terms form an arithmetic progression, the last three terms form a geometric progression, and the first and fourth terms differ by $30.$ Find the sum of the four terms.


 
 \textbf{Answer: }129



\textbf{Solution}

Denote the first term as $a$, and the common difference between the first three terms as $d$. The four numbers thus are in the form $a,\ a+d,\ a+2d,\ \frac{(a + 2d)^2}{a + d}$. 

 Since the first and fourth terms differ by $30$, we have that $\frac{(a + 2d)^2}{a + d} - a = 30$. Multiplying out by the denominator, 
 \[(a^2 + 4ad + 4d^2) - a(a + d) = 30(a + d).\] This simplifies to $3ad + 4d^2 = 30a + 30d$, which upon rearranging yields $2d(2d - 15) = 3a(10 - d)$.

 Both $a$ and $d$ are positive integers, so $2d - 15$ and $10 - d$ must have the same sign. Try if they are both \href{https://artofproblemsolving.com/wiki/index.php?title=Positive}{positive} (notice if they are both \href{https://artofproblemsolving.com/wiki/index.php?title=Negative}{negative}, then $d > 10$ and $d < \frac{15}{2}$, which is a contradiction). Then, $d = 8, 9$. Directly substituting and testing shows that $d \neq 8$, but that if 

 


	\item (\textit{AIME I 2003}) An integer between $1000$ and $9999,$ inclusive, is called balanced if the sum of its two leftmost digits equals the sum of its two rightmost digits. How many balanced integers are there?


 
 \textbf{Answer: }615



\textbf{Solution 1}

If the common sum of the first two and last two digits is $n$, such that $1 \leq n \leq 9$, there are $n$ choices for the first two digits and $n + 1$ choices for the second two digits (since zero may not be the first digit).  This gives $\sum_{n = 1}^9 n(n + 1) = 330$ balanced numbers.  If the common sum of the first two and last two digits is $n$, such that $10 \leq n \leq 18$, there are $19 - n$ choices for both pairs.  This gives $\sum_{n = 10}^{18} (19 - n)^2 = \sum_{n = 1}^9 n^2 = 285$ balanced numbers.  Thus, there are in total $330 + 285 = \boxed{615}$ balanced numbers.

 Both summations may be calculated using the formula for the \href{https://artofproblemsolving.com/wiki/index.php?title=Perfect\_square}{sum of consecutive squares}, namely $\sum_{k=1}^n k^2 = \frac{n(n+1)(2n+1)}{6}$.

 



\textbf{Solution 2 (Painful Casework)}

Call the number $\overline{abcd}$. Then $a+b=c+d$. Set $a+b=x$.

 Clearly, $0\le x \le18$.

 If $x=0$: $0000$ is not acceptable.

 If $x=1$: The only case is $1001$ or $1010$. 2 choices.

 If $x=2$: then since $a\neq0$, $a=1=b$ or $a=2, b=0$. There are 3 choices for $(c,d)$: $(2,0), (0, 2), (1, 1)$. $2*3=6$ here.

 If $x=3$: Clearly, $a\neq b$ because if so, the sum will be even, not odd. Counting $(a,b)=(3,0)$, we have $4$ choices. Subtracting that, we have $3$ choices. Since it doesn't matter whether $c=0$ or $d=0$, we have 4 choices for $(c,d)$. So $3*4=12$ here.

 If $x=4$: Continue as above. $4$ choices for $(a,b)$. $5$ choices for $(c,d)$. $4*5=20$ here.

 If $x=5$: You get the point. $5*6=30$.

 If $x=6$: $6*7=42$.

 If $x=7$: $7*8=56$.

 If $x=8$: $8*9=72$.

 If $x=9$: $9*10=90$.

 Now we need to be careful because if $x=10$, $(c,d)=(0,10)$ is not valid. However, we don't have to worry about $a\neq0$.

 If $x=10$: $(a,b)=(1,9), (2, 8), ..., (9, 1)$. Same thing for $(c,d)$. $9*9=81$.

 If $x=11$: We start at $(a,b)= (2,9)$. So $8*8$.

 Continue this pattern until $x=18: 1*1=1$. Add everything up: we have $\boxed{615}$.

 



\textbf{Solution 3}

We ignore the requirement that the first digit is non-zero, and do casework on the sum of the sum of the pairs of digits.

 If two digits $a$ and $b$ sum to $0$, we have $1$ possibility: $(a,b) = (0,0)$

 If $a+b = 1$, we have $2$ possibilities: $(a,b) = (0,1)$ and $(a,b) = (1,0)$

 $a+b = 2$: $(0,2)$, $(2,0)$, and $(1,1)$ are the only $3$ possibilities.

 We notice a pattern: for all $k \leq 9$, there are $k+1$ ordered pairs of digits $(a,b)$ such that $a+b = k$. Then, testing for $10 \leq k \leq 18$, we notice that there are $(18 - k) + 1$ ordered pairs with a sum of $k$.

 So the number of ways to pick $4$ digits satisfying the constraint is 
 \[(1)(1) + (2)(2) \dots (10)(10) + (9)(9) + \dots + (1)(1) = \frac{(10)(11)(21)}{6} + \frac{(9)(10)(19)}{6} = 670\]

 We have to subtract out the cases where the first digit is $0$, which is $1+2+\dots+10  = 55$

 So our answer is $670-55 = \boxed{615}$

 This solution takes some inspiration from dice. In a way, we are counting the number of ways to roll a given number with a pair of $10$-sided dice. 

 


	\item (\textit{AIME I 2003}) Triangle $ABC$ is isosceles with $AC = BC$ and $\angle ACB = 106^\circ.$ Point $M$ is in the interior of the triangle so that $\angle MAC = 7^\circ$ and $\angle MCA = 23^\circ.$ Find the number of degrees in $\angle CMB.$


 
 \textbf{Answer: }083



\textbf{Solution 1}


                    \begin{figure}[H]
                        \centering
                        \includegraphics[width=0.5\textwidth]{Solution/AIME_I_2003_10_Solution_Figure_1.png}
                    \end{figure}
                    

 Take point $N$ inside $\triangle ABC$ such that $\angle CBN = 7^\circ$ and $\angle BCN = 23^\circ$.

 $\angle MCN = 106^\circ - 2\cdot 23^\circ = 60^\circ$. Also, since $\triangle AMC$ and $\triangle BNC$ are congruent (by ASA), $CM = CN$. Hence $\triangle CMN$ is an \href{https://artofproblemsolving.com/wiki/index.php?title=Equilateral\_triangle}{equilateral triangle}, so $\angle CNM = 60^\circ$.

 Then $\angle MNB = 360^\circ - \angle CNM - \angle CNB = 360^\circ - 60^\circ - 150^\circ = 150^\circ$. We now see that $\triangle MNB$ and $\triangle CNB$ are congruent. Therefore, $CB = MB$, so $\angle CMB = \angle MCB = \boxed{83^\circ}$.

 



\textbf{Solution 2}

From the givens, we have the following \href{https://artofproblemsolving.com/wiki/index.php?title=Angle}{angle} \href{https://artofproblemsolving.com/wiki/index.php?title=Measure}{measures}: $m\angle AMC = 150^\circ$, $m\angle MCB = 83^\circ$.  If we define $m\angle CMB = \theta$ then we also have $m\angle CBM = 97^\circ - \theta$.  Then apply the \href{https://artofproblemsolving.com/wiki/index.php?title=Law\_of\_Sines}{Law of Sines} to triangles $\triangle AMC$ and $\triangle BMC$ to get

 
 \[\frac{\sin 150^\circ}{\sin 7^\circ} = \frac{AC}{CM} = \frac{BC}{CM} = \frac{\sin \theta}{\sin (97^\circ - \theta)}\]

 Clearing \href{https://artofproblemsolving.com/wiki/index.php?title=Denominator}{denominators}, evaluating $\sin 150^\circ = \frac 12$ and applying one of our \href{https://artofproblemsolving.com/wiki/index.php?title=Trigonometric\_identities}{trigonometric identities} to the result gives

 
 \[\frac{1}{2} \cos (7^\circ - \theta )= \sin 7^\circ \sin \theta\]

 and multiplying through by 2 and applying the \href{https://artofproblemsolving.com/wiki/index.php/Trigonometric\_identities\#Double-angle\_identities\%7C}{double angle formulas} gives

 
 \[\cos 7^\circ\cos\theta + \sin7^\circ\sin\theta = 2 \sin7^\circ \sin\theta\]

 and so $\cos 7^\circ \cos \theta = \sin 7^\circ \sin\theta \Longleftrightarrow \tan 7^{\circ} = \cot \theta$; since $0^\circ < \theta < 180^\circ$, we must have $\theta = 83^\circ$, so the answer is $\boxed{83}$.

 



\textbf{Solution 3}

\href{https://artofproblemsolving.com/wiki/index.php?title=Without\_loss\_of\_generality}{Without loss of generality}, let $AC = BC = 1$.  Then, using the \href{https://artofproblemsolving.com/wiki/index.php?title=Law\_of\_Sines}{Law of Sines} in triangle $AMC$, we get $\frac {1}{\sin 150} = \frac {MC}{\sin 7}$, and using the sine addition formula to evaluate $\sin 150 = \sin (90 + 60)$, we get $MC = 2 \sin 7$.  

 Then, using the \href{https://artofproblemsolving.com/wiki/index.php?title=Law\_of\_Cosines}{Law of Cosines} in triangle $MCB$, we get $MB^2 = 4\sin^2 7 + 1 - 4\sin 7(\cos 83) = 1$, since $\cos 83 = \sin 7$. So triangle $MCB$ is isosceles, and $\angle CMB = \boxed{83}$.

 



\textbf{Solution 4}

Note: A diagram would be much appreciated; I cannot make one since I'm bad at asymptote. Also, please make this less cluttered :) ~tauros

 First, take point $E$ outside of $\triangle{ABC}$ so that $\triangle{CEB}$ is equilateral. Then, connect $A$, $C$, and $M$ to $E$. Also, let $ME$ intersect $AB$ at $F$. $\angle{MCE} = 83^\circ - 60^\circ = 23^\circ$, $CE = AB$, and (trivially) $CM = CM$, so $\triangle{MCE} \cong  \triangle{MCA}$ by SAS congruence. Also, $\angle{CMA} = \angle{CME} = 150^\circ$, so $\angle{AME} = 60^\circ$, and $AM = ME$,making $\triangle{AME}$ also equilateral. (it is isosceles with a $60^\circ$ angle) $\triangle{MAF} \cong \triangle{EAF}$ by SAS ($MA = AE$, $AF = AF$, and $m\angle{MAF} = m\angle{EAF} = 30^\circ$), and $\triangle{MAB} \cong \triangle{EAB}$ by SAS ($MA = AE$, $AB = AB$, and $m\angle{MAB} = m\angle{EAB} = 30^\circ$). Thus, $\triangle{BME}$ is isosceles, with $m\angle{BME} = m\angle{BEM} = 60^\circ + 7^\circ = 67^\circ$. Also, $\angle{EMB} + \angle {CMB} = \angle{CME} = 150^\circ$, so $\angle{CME} = 150^\circ - 67^\circ  = \boxed{83^\circ}$.

 



\textbf{Solution 5 (Ceva)}

Noticing that we have three concurrent cevians, we apply Ceva's theorem:

 
 \[(\sin \angle ACM)(\sin \angle BAM)(\sin \angle CBM) = (\sin \angle CAM)(\sin \angle ABM)(\sin \angle BCM)\]
 \[(\sin 23)(\sin 30)(\sin x) = (\sin 7)(\sin 37-x)(\sin 83)\]

 using the fact that $\sin 83 = \cos 7$ and $(\sin 7)(\cos 7) = 1/2 (\sin 14)$ we have:

 
 \[(\sin 23)(\sin x) = (\sin 14)(\sin 37-x)\]

 By inspection, $x=14^\circ$ works, so the answer is $180-83-14= \boxed{083}$

 



\textbf{Solution 6}

Let $\angle{APC} = \theta^{\circ}$Using sine rule on $\triangle{APB}, \triangle{APC}$, letting $AP=d$ we get :$\frac{d}{1} = \frac{\sin{7^{\circ}}}{\sin{150^{\circ}}} = 2\sin{7^{\circ}}= \frac{\sin{14^{\circ}}}{\cos{7^{\circ}}}= \frac{\sin{14^{\circ}}}{\sin{83^{\circ}}}= \frac{\sin{(97-\theta)^{\circ}}}{\sin{\theta^{\circ}}}$Simplifying, we get that$\cos{(14-\theta)^{\circ}}-\cos{(14+\theta)^{\circ}}=\cos{(14-\theta)^{\circ}}-\cos{(180-\theta)^{\circ}},$ from where $\cos{(14-\theta)^{\circ}}=\cos{(180-\theta)^{\circ}}$Simplifying more, we get that $\sin{97^{\circ}} \cdot \sin{(\theta-83)^{\circ}} = 0$, so $\theta = 83^{\circ}$NOTE: The simplifications were carried out by the product-to-sum and sum-to-product identities ~Prabh1512

 



\textbf{Solution 7 (Thorough & Has diagram)}


                    \begin{figure}[H]
                        \centering
                        \includegraphics[width=0.5\textwidth]{Solution/AIME_I_2003_10_Solution_Figure_2.png}
                    \end{figure}
                    Construct $\overline{CP} \perp \overline{AB}$. Then, since $\triangle{ABC}$ is isosceles, $\overline{CP}$ is the perpindicular bisector of $\overline{AB}$

 Then, reflect $M$ over $\overline{CP}$ to $M'$. This also maps $A\rightarrow B$ and therefore $\angle{CMB}\cong\angle{CM'A}$.

 Observe that $\angle{CMA}=180-\angle{ACM}-\angle{MAC}=150^\circ$. Further, since our reflection maps $A\rightarrow A$ and $M\rightarrow M'$ $\angle{M'CB}=\angle{ACM}=23^\circ \therefore \angle{M'CM}=60^\circ$. Since C is on the line of refelction, $CM\cong CM'$ and therefore $\triangle{CMM'}$ is equilateral. As such, $\angle{M'MA}=150^\circ$. Then, if $\angle{M'AM}=\theta, \angle{AM'M}=30-\theta$.

 By law of sines on $\triangle{CMA}$, 
 \[\frac{MA}{\sin(23)}=\frac{CM}{\sin(7)}\implies MA = \frac{CM\sin(23)}{\sin(7)}\]

 By law of sines on $\triangle{AMM'}$, 
 \[\frac{MA}{\sin(30-\theta)}=\frac{MM'}{\sin(\theta)}\implies MA = \frac{MM'\sin(30-\theta)}{\sin(\theta)}\]So:
 \[\frac{MM'\sin(30-\theta)}{\sin(\theta)}=\frac{CM\sin(23)}{\sin(7)}\implies\frac{\sin(30-\theta)}{\sin(\theta)}=\frac{\sin(23)}{\sin(7)}\because \triangle{CMM'} \text{is equilateral}\]This obviously gives $\theta = 7, 30-\theta = 23$

 Finally, 
 \[\angle{CMB}=\angle{CM'A}=\angle{CM'M}+\angle{MM'A}=60+23=\boxed{083}\]~Stereotypicalmathnerd

 


	\item (\textit{AIME I 2003}) An angle $x$ is chosen at random from the interval $0^\circ < x < 90^\circ.$ Let $p$ be the probability that the numbers $\sin^2 x, \cos^2 x,$ and $\sin x \cos x$ are not the lengths of the sides of a triangle. Given that $p = d/n,$ where $d$ is the number of degrees in $\arctan m$ and $m$ and $n$ are positive integers with $m + n < 1000,$ find $m + n.$


 
 \textbf{Answer: }092



\textbf{Solution}

Note that the three expressions are symmetric with respect to interchanging $\sin$ and $\cos$, and so the probability is symmetric around $45^\circ$.  Thus, take $0 < x < 45$ so that $\sin x < \cos x$.  Then $\cos^2 x$ is the largest of the three given expressions and those three lengths not forming a \href{https://artofproblemsolving.com/wiki/index.php?title=Triangle}{triangle} is equivalent to a violation of the \href{https://artofproblemsolving.com/wiki/index.php?title=Triangle\_inequality}{triangle inequality}

 
 \[\cos^2 x > \sin^2 x + \sin x \cos x\]

 This is equivalent to 

 
 \[\cos^2 x - \sin^2 x > \sin x \cos x\]

 and, using some of our \href{https://artofproblemsolving.com/wiki/index.php?title=Trigonometric\_identities}{trigonometric identities}, we can re-write this as $\cos 2x > \frac 12 \sin 2x$. Since we've chosen $x \in (0, 45)$, $\cos 2x > 0$ so 

 
 \[2 > \tan 2x \Longrightarrow  x < \frac 12 \arctan 2.\]

 The \href{https://artofproblemsolving.com/wiki/index.php?title=Probability}{probability} that $x$ lies in this range is $\frac 1{45} \cdot \left(\frac 12 \arctan 2\right) = \frac{\arctan 2}{90}$ so that $m = 2$, $n = 90$ and our answer is $\boxed{092}$.

 



\textbf{Solution 2 (Complementary Counting)}

We seek a complementary counting argument, where we look for the probability that $\sin^2 x$, $\cos^2 x$ and $\sin x \cos x$ form the side lengths of a triangle.

 By the triangle inequality, we must have the following three inequalities to be true: 
 \[\sin^2 x + \cos^2 x > \sin x \cos x\] 
 \[\sin^2 x + \sin x \cos x > \cos^2 x\] 
 \[\cos^2 x + \sin x \cos x > \sin^2 x\]The first inequality will always hold since we have $\sin^2 x + \cos^2 x = 1$, and $1 > \sin x \cos x$ for all $x$ (The maximum value of $\sin x \cos x$ is $\frac{1}{2}$ when $\sin x = \cos x = \frac{\sqrt{2}}{2}$).

 Now, we examine the second inequality $\sin^2 x + \sin x \cos x > \cos^2 x$. If we subtract $\sin^2 x$ from both sides, we have $\sin x \cos x > \cos^2 x - \sin^2 x$. Aha! This resembles our sine and cosine double angle identities. Therefore, our inequality is now $\sin 2x > 2\cos 2x$. We can divide both sides by $\cos 2x$ and we have $\tan 2x > 2$.  The solutions to this occur when $45 \geq x > \frac{\arctan 2}{2}$. 

 (To understand why it must be $x >$, we can draw the unit circle, and notice as x moves from $\frac{\arctan 2}{2}$ to $90$, $\tan x$ approaches $\infty$. We must cap $x$ at $45$, since if $x > 45$, $2x > 90$, and $\tan x$ will be negative.)

 
Next, we examine the third inequality, $\cos^2 x + \sin x \cos x > \sin^2 x$. Once again, we can get our double angle identities for sine and cosine by subtracting $\cos^2 2x$ from both sides. We have, $\sin x \cos x > \sin^2 x -\cos^2 x \to \sin 2x > -2\cos 2x$. 

 Next, we again, divide by $\cos 2x$ to produce a $\tan 2x$ (we do this because one trig function is easier to deal with than 2). However, if $\cos 2x > 0$, we do not need to flip the sign since $\sin 2x >0$, and so if $\cos 2x >0$, all values for which that is true satisfy the inequality. So we only consider if $\cos 2x < 0$, and when we divide by a negative, we must flip the sign. Thus we have $\tan 2x < -2$. 

 We can take the $\arctan$ of both sides, and we have $\frac{\arctan -2}{2}> x \geq 45$. Once again, to better understand this, we can draw the angle $x$ for which $\tan 2x = -2$, and we notice as $2x$ moves to $x=90$, $\tan 2x$ approaches $- \infty$. We must cap $x$ at $45$ since if $x<45$, we have $\tan 2x > 0$. 

 Notice that if we draw the terminal points for $\frac{\arctan 2}{2}$ and $\frac{\arctan -2}{2}$, they have the same smaller angle with the x and y axis respectively. This means the range of degree measures for which our inequalities hold is $90 - \frac{\arctan 2}{2} > x > \frac{\arctan 2}{2}$ which has an area of $90 - \arctan 2$. However, we want the complement of this, which has an area of $90 - (90 - \arctan 2) = \arctan 2$. Therefore, the desired probability is $\frac{\arctan2}{90}$, and so $m+n=2+90=92$.

 -BossLu99

 


 


 


	\item (\textit{AIME I 2003}) In convex quadrilateral $ABCD, \angle A \cong \angle C, AB = CD = 180,$ and $AD \neq BC.$ The perimeter of $ABCD$ is 640. Find $\lfloor 1000 \cos A \rfloor.$ (The notation $\lfloor x \rfloor$ means the greatest integer that is less than or equal to $x.$)


 
 \textbf{Answer: }777



\textbf{Solution 1}


                    \begin{figure}[H]
                        \centering
                        \includegraphics[width=0.5\textwidth]{Solution/AIME_I_2003_12_Solution_Figure_1.png}
                    \end{figure}
                    

 By the \href{https://artofproblemsolving.com/wiki/index.php?title=Law\_of\_Cosines}{Law of Cosines} on $\triangle ABD$ at angle $A$ and on $\triangle BCD$ at angle $C$ (note $\angle C = \angle A$),

 
 \[180^2 + AD^2 - 360 \cdot AD \cos A = 180^2 + BC^2 - 360 \cdot BC \cos A\]
 \[(AD^2 - BC^2) = 360(AD - BC) \cos A\]
 \[(AD - BC)(AD + BC) = 360(AD - BC) \cos A\]
 \[(AD + BC) = 360 \cos A\]We know that $AD + BC = 640 - 360 = 280$. $\cos A = \dfrac{280}{360} =  \dfrac{7}{9} = 0.777 \ldots$

 $\lfloor 1000 \cos  A \rfloor = \boxed{777}$.

 



\textbf{Solution 2}

Notice that $AB = CD$, and $BD = DB$, and $\angle{DAB} \cong \angle{BCD}$, so we have side-side-angle matching on triangles $ABD$ and $CDB$.  Since the problem does not allow $\triangle{ABD} \cong \triangle{CDB}$, we know that $\angle{ADB}$ is not a right angle, and there is a unique other triangle with the matching side-side-angle.

 
                    \begin{figure}[H]
                        \centering
                        \includegraphics[width=0.5\textwidth]{Solution/AIME_I_2003_12_Solution_Figure_2.png}
                    \end{figure}
                    

 Overlay the triangles $\triangle BDA$ and $\triangle BDC$ on each other as in the diagram above (where the red labels correspond to $\triangle BDA$ and the blue labels correspond to $\triangle BDC$). Here we assume without loss of generality that $BC>AD$. Furthermore, let $\theta$ be the angle $A$ referenced in the problem; we need to find $\lfloor 1000\cos\theta\rfloor$.

 Since the perimeter of $ABCD$ is $640$, we have $AD+BC=640-360=280$. Thus let $BC=a$ and $AD=280-a$ for some positive real number $a$. But the sides that correspond to $\overline{BD}$ above are congruent, so we can drop a perpendicular from the topmost point to a point $E$, where the base of the isosceles triangle is bisected. Notice that the base of the isosceles triangle has length $a-(280-a)=2a-280$, so in the diagram above, $AE=280-a+\frac{1}{2}(2a-280)=280-a+a-140=140$.

 Looking above at the right triangle containing $\theta$, we see that $\cos\theta=\frac{AE}{180}=\frac{140}{180}=\frac{7}{9}$. Hence $\left\lfloor 1000\cdot\frac{7}{9}\right\rfloor=\boxed{777}$ is our answer.

 



\textbf{Solution 3}

Start the same as solution 1. We get 
 \[180^2+a^2-360a \cos A = 180^2+b^2-360b \cos A \Rightarrow a^2-360a \cos A = b^2-360b \cos A,\]where $a$ is the length of $BC$ and $b$ is the length of $AD$. Let the common value of $a^2-360a \cos A$ and $b^2-360b \cos A$ be $c$. Then, the quadratic in $x$ 
 \[x^2-360 \cos A \cdot x - c = 0\]has solutions $a$ and $b$. Therefore, by Vieta's, $a+b = 360 \cos A$. However, we know that the perimeter of $ABCD$ is $640$, so $a+b+180+180=640$, so $a+b=280$. Therefore, 
 \[360 \cos A = 280 \Rightarrow \cos A = \frac{280}{360} \Rightarrow \cos A = \frac{7}{9} \Rightarrow \lfloor 1000 \cos A \rfloor = \boxed{777}.\]

 


	\item (\textit{AIME I 2003}) Let $N$ be the number of positive integers that are less than or equal to 2003 and whose base-2 representation has more 1's than 0's. Find the remainder when $N$ is divided by 1000.


 
 \textbf{Answer: }155



\textbf{Solution 1}

In base-$2$ representation, all positive numbers have a leftmost digit of $1$. Thus there are ${n \choose k}$ numbers that have $n+1$ digits in base $2$ notation, with $k+1$ of the digits being $1$'s. 

 In order for there to be more $1$'s than $0$'s, we must have $k+1 > \frac{n+1}{2} \implies k > \frac{n-1}{2} \implies k \ge \frac{n}{2}$. Therefore, the number of such numbers corresponds to the sum of all numbers on or to the right of the vertical line of symmetry in \href{https://artofproblemsolving.com/wiki/index.php?title=Pascal\%27s\_Triangle}{Pascal's Triangle}, from rows $0$ to $10$ (as $2003 < 2^{11}-1$). Since the sum of the elements of the $r$th row is $2^r$, it follows that the sum of all elements in rows $0$ through $10$ is $2^0 + 2^1 + \cdots + 2^{10} = 2^{11}-1 = 2047$. The center elements are in the form ${2i \choose i}$, so the sum of these elements is $\sum_{i=0}^{5} {2i \choose i} = 1 + 2 +6 + 20 + 70 + 252 = 351$.

 The sum of the elements on or to the right of the line of symmetry is thus $\frac{2047 + 351}{2} = 1199$. However, we also counted the $44$ numbers from $2004$ to $2^{11}-1 = 2047$. Indeed, all of these numbers have at least $6$ $1$'s in their base-$2$ representation, as all of them are greater than $1984 = 11111000000_2$, which has $5$ $1$'s. Therefore, our answer is $1199 - 44 = 1155$, and the remainder is $\boxed{155}$.

 



\textbf{Solution 2}

We seek the number of allowed numbers which have $k$ 1's, not including the leading 1, for $k=0, 1, 2, \ldots , 10$.

 For $k=0,\ldots , 4$, this number is

 $\binom{k}{k}+\binom{k+1}{k}+\cdots+\binom{2k}{k}$.

 By the Hockey Stick Identity, this is equal to $\binom{2k+1}{k+1}$. So we get

 $\binom{1}{1}+\binom{3}{2}+\binom{5}{3}+\binom{7}{4}+\binom{9}{5}=175$.

 For $k=5,\ldots , 10$, we end on $\binom{10}{k}$ - we don't want to consider numbers with more than 11 digits. So for each $k$ we get

 $\binom{k}{k}+\binom{k+1}{k}+\ldots+\binom{10}{k}=\binom{11}{k+1}$

 again by the Hockey Stick Identity. So we get

 $\binom{11}{6}+\binom{11}{7}+\binom{11}{8}+\binom{11}{9}+\binom{11}{10}+\binom{11}{11}=\frac{2^{11}}{2}=2^{10}=1024$.

 The total is $1024+175=1199$. Subtracting out the $44$ numbers between $2003$ and $2048$ gives $1155$. Thus the answer is $155$.

 



\textbf{Solution 3}

We will count the number of it $< 2^{11}=2048$ instead of $2003$ (In other words, the length of the base-2 representation is at most $11$. If there are even digits, $2n$, then the leftmost digit is $1$, the rest, $2n-1$, has odd number of digits. In order for the base-2 representation to have more $1$'s, we will need more $1$ in the remaining $2n-1$ than $0$'s. Using symmetry, this is equal to$\frac{2^9+2^7+..+2^1}{2}$Using similar argument where there are odd amount of digits. The remaining even amount of digit must contains the number of $1$'s at least as the number of $0$'s. So it's equal to$\frac{\binom{10}{5}+2^{10}+\binom{8}{4}+2^8+\binom{6}{3}+2^6+...+\binom{0}{0}+2^0}{2}$Summing both cases, we have $\frac{2^0+2^1+..+2^{10}+\binom{0}{0}+..+\binom{10}{5}}{2} = 199$. There are $44$ numbers between $2004$ and $2047$ inclusive that satisfy it. So the answer is $199-44=\boxed{155}$

 


	\item (\textit{AIME I 2003}) The decimal representation of $m/n,$ where $m$ and $n$ are relatively prime positive integers and $m < n,$ contains the digits 2, 5, and 1 consecutively, and in that order. Find the smallest value of $n$ for which this is possible.


 
 \textbf{Answer: }127



\textbf{Solution}

To find the smallest value of $n$, we consider when the first three digits after the decimal point are $0.251\ldots$. 

 
Otherwise, suppose the number is in the form of $\frac{m}{n} = 0.X251 \ldots$, where $X$ is a string of $k$ digits and $n$ is small as possible. Then $10^k \cdot \frac{m}{n} - X = \frac{10^k m - nX}{n} = 0.251 \ldots$. Since $10^k m - nX$ is an integer and $\frac{10^k m - nX}{n}$ is a fraction between $0$ and $1$, we can rewrite this as $\frac{10^k m - nX}{n} = \frac{p}{q}$, where $q \le n$. Then the fraction $\frac pq = 0.251 \ldots$ suffices. 

 Thus we have $\frac{m'}{n} = 0.251\ldots$, or 

 $\frac{251}{1000} \le \frac{m'}{n} < \frac{252}{1000} \Longleftrightarrow 251n \le 1000m' < 252n \Longleftrightarrow n \le 250(4m'-n) < 2n.$

 As $4m' > n$, we know that the minimum value of $4m' - n$ is $1$; hence we need $250 < 2n \Longrightarrow 125 < n$. Since $4m' - n = 1$, we need $n + 1$ to be divisible by $4$, and this first occurs when $n = \boxed{ 127 }$ (note that if $4m'-n > 1$, then $n > 250$). Indeed, this gives $m' = 32$ and the fraction $\frac {32}{127}\approx 0.25196 \ldots$).

 



\textbf{Solution 2}

Rewrite the problem as having the smallest $n$ such that we can find an positive integer $m$ such that $0<\frac{m}{n}-\frac{251}{1000}<\frac{1}{1000}$.

 We can rewrite the expression as $\frac{1000m-251n}{1000n}$, and we need $251n+x$ (where $x$ is the difference in the fraction, and ranging from $(1,2,...n-1)$ to be $0$ mod $1000$. We see that $n$ must be $3$ mod $4$ to have this happen (as this reduces the distance between the expression and $1000$.

 Rewriting $n$ as $4k+3$, we get that $251(4k+3)+(4k+2)$ turns into $8k+755$, and this has to be greater than or equal to $1000$. The least $k$ that satisfies this is $31$, and we consequently get that the least value of $n$ is $127$. -dragoon -minor edits by Mathkiddie

 


	\item (\textit{AIME I 2003}) In $\triangle ABC, AB = 360, BC = 507,$ and $CA = 780.$ Let $M$ be the midpoint of $\overline{CA},$ and let $D$ be the point on $\overline{CA}$ such that $\overline{BD}$ bisects angle $ABC.$ Let $F$ be the point on $\overline{BC}$ such that $\overline{DF} \perp \overline{BD}.$ Suppose that $\overline{DF}$ meets $\overline{BM}$ at $E.$ The ratio $DE: EF$ can be written in the form $m/n,$ where $m$ and $n$ are relatively prime positive integers. Find $m + n.$


 
 \textbf{Answer: }289



\textbf{Solution 1}

In the following, let the name of a point represent the mass located there. Since we are looking for a ratio, we assume that $AB=120$, $BC=169$, and $CA=260$ in order to simplify our computations.

 First, reflect point $F$ over angle bisector $BD$ to a point $F'$.

 
                    \begin{figure}[H]
                        \centering
                        \includegraphics[width=0.5\textwidth]{Solution/AIME_I_2003_15_Solution_Figure_1.png}
                    \end{figure}
                    

 As $BD$ is an angle bisector of both triangles $BAC$ and $BF'F$, we know that $F'$ lies on $AB$. We can now balance triangle $BF'C$ at point $D$ using mass points.

 By the \href{https://artofproblemsolving.com/wiki/index.php?title=Angle\_Bisector\_Theorem}{Angle Bisector Theorem}, we can place \href{https://artofproblemsolving.com/wiki/index.php?title=Mass\_points}{mass points} on $C,D,A$ of $120,\,289,\,169$ respectively. Thus, a mass of $\frac {289}{2}$ belongs at both $F$ and $F'$ because BD is a median of triangle $BF'F$ . Therefore, $CB/FB=\frac{289}{240}$.

 Now, we reassign mass points to determine $FE/FD$. This setup involves $\triangle CFD$ and \href{https://artofproblemsolving.com/wiki/index.php?title=Transversal}{transversal} $MEB$. For simplicity, put masses of $240$ and $289$ at $C$ and $F$ respectively. To find the mass we should put at $D$, we compute $CM/MD$. Applying the Angle Bisector Theorem again and using the fact $M$ is a midpoint of $AC$, we find
 \[\frac {MD}{CM} = \frac {\frac{169}{289}\cdot 260 - 130}{130} = \frac {49}{289}\]At this point we could find the mass at $D$ but it's unnecessary.
 \[\frac {DE}{EF} = \frac {F}{D} = \frac {F}{C}\cdot\frac {C}{D} = \frac {289}{240}\cdot\frac {49}{289} = \boxed{\frac {49}{240}}\]and the answer is $49 + 240 = \boxed{289}$.

 



\textbf{Solution 2}

By the Angle Bisector Theorem, we know that $[CBD]=\frac{169}{289}[ABC]$. Therefore, by finding the area of triangle $CBD$, we see that 
 \[\frac{507\cdot BD}{2}\sin\frac{B}{2}=\frac{169}{289}[ABC].\]Solving for $BD$ yields 
 \[BD=\frac{2[ABC]}{3\cdot289\sin\frac{B}{2}}.\]Furthermore, $\cos\frac{B}{2}=\frac{BD}{BF}$, so 
 \[BF=\frac{BD}{\cos\frac{B}{2}}=\frac{2[ABC]}{3\cdot289\sin\frac{B}{2}\cos\frac{B}{2}}.\]Now by the identity $2\sin\frac{B}{2}\cos\frac{B}{2}=\sin B$, we get 
 \[BF=\frac{4[ABC]}{3\cdot289\sin B}.\]But then $[ABC]=\frac{360\cdot 507}{2}\sin B$, so $BF=\frac{240}{289}\cdot 507$. Thus $BF:FC=240:49$.

 Now by the Angle Bisector Theorem, $CD=\frac{169}{289}\cdot 780$, and we know that $MC=\frac{1}{2}\cdot 780$ so $DM:MC=\frac{169}{289}-\frac{1}{2}:\frac{1}{2}=49:289$.

 We can now use mass points on triangle CBD. Assign a mass of $240\cdot 49$ to point $C$. Then $D$ must have mass $240\cdot 289$ and $B$ must have mass $49\cdot 49$. This gives $F$ a mass of $240\cdot 49+49\cdot 49=289\cdot 49$. Therefore, $DE:EF=\frac{289\cdot 49}{240\cdot 289}=\frac{49}{240}$, giving us an answer of $\boxed{289}.$

 



\textbf{Solution 3}

Let $\angle{DBM}=\theta$ and $\angle{DBC}=\alpha$. Then because $BM$ is a median we have $360\sin{(\alpha+\theta)}=507\sin{(\alpha-\theta)}$. Now we know 
 \[\sin{(\alpha+\theta)}=\sin{\alpha}\cos{\theta}+\sin{\theta}\cos{\alpha}=\dfrac{DF\cdot BD}{BF\cdot BE}+\dfrac{DE\cdot BD}{BE\cdot BF}=\dfrac{BD(DF+DE)}{BF\cdot BE}\] Expressing the area of $\triangle{BEF}$ in two ways we have 
 \[\dfrac{1}{2}BE\cdot BF\sin{(\alpha-\theta)}=\dfrac{1}{2}EF\cdot BD\] so 
 \[\sin{(\alpha-\theta)}=\dfrac{EF\cdot BD}{BF\cdot BE}\] Plugging this in we have 
 \[\dfrac{360\cdot BD(DF+DE)}{BF\cdot BE}=\dfrac{507\cdot BD\cdot EF}{BF\cdot BE}\] so $\dfrac{DF+DE}{EF}=\dfrac{507}{360}$. But $DF=DE+EF$, so this simplifies to $1+\dfrac{2DE}{EF}=\dfrac{507}{360}=\dfrac{169}{120}$, and thus $\dfrac{DE}{EF}=\dfrac{49}{240}$, and $m+n=\boxed{289}$.

 



\textbf{Solution 4 (Overpowered Projective Geometry!!)}

Firstly, angle bisector theorem yields $\frac{CD}{AD} = \frac{507}{360} = \frac{169}{120}$. We're given that $AM=MC$. Therefore, the cross ratio

 
 \[(A,C;M,D) = \frac{AM(CD)}{AD(MC)} = \frac{169}{120}\]

 We need a fourth point for this cross ratio to be useful, so reflect point $F$ over angle bisector $BD$ to a point $F'$. Then $\triangle BFF'$ is isosceles and $BD$ is an altitude so $DF = DF'$. Therefore, 

 
 \[(A,C;M,D) = (F,F';D,E) \implies \frac{FD(EF')}{EF(DF')} = \frac{EF'}{EF} = \frac{169}{120}\]

 All that's left is to fiddle around with the ratios:

 
 \[\frac{EF'}{EF} = \frac{ED+DF'}{EF} = \frac{EF+2DE}{EF} = 1\ +\ 2\left(\frac{DE}{EF}\right) \implies \frac{DE}{EF} = \frac{49}{240} \implies \boxed{289}\]

 



\textbf{Solution 5 (Menelaus + Mass Points)}

Extend $DF$ to intersect with the extension of $AB$ at $G$. Notice that $\triangle{BDF} \cong \triangle{BDG}$, so $GD=DF$. We now use Menelaus on $\triangle{GBF}$, as $A$, $D$, and $C$ are collinear; this gives us $\frac{GA}{BA} \cdot \frac{BC}{FC} \cdot \frac{DF}{GD}=1$. As $GD=DF$, we have $\frac{GA}{AB}=\frac{FC}{BC}$, hence $\frac{GA}{120}=\frac{FC}{169}$. Reflect $G$ over $A$ to $G'$. Note that $\frac{G'A}{BA}=\frac{FC}{BC}$, and reflexivity, hence $\triangle{ABC} \sim \triangle{BG'F}$. It's easily concluded from this that $G'F \parallel AC$, hence $G'F \parallel AD$. As $GD=DF$, we have $AD$ is a midsegment of $\triangle{GG'F}$, thus $G'F = 2AD$. We now focus on the ratio $\frac{BF}{BC}$. From similarity, we have $\frac{BF}{BC}=\frac{G'F}{AC}=\frac{2AD}{AC}$. By the angle bisector theorem, we have $AD:DC=120:169$, hence $AD:AC=120:289$, so $\frac{BF}{BC}=\frac{240}{289}$. We now work out the ratio $\frac{DM}{MC}$. $\frac{DM}{MC}=\frac{CD-MC}{MC}=\frac{CD}{MC}-1=\frac{2CD}{AC}-1=\frac{338}{289}-1=\frac{49}{289}$. We now use mass points on $\triangle{BDC}$. We let the mass of $C$ be $240\cdot 49$, so the mass of $B$ is $49 \cdot 49$ and the mass of $D$ is $289\cdot 240$. Hence, the mass of $F$ is $289\cdot 49$, so the ratio $\frac{DE}{EF}=\frac{49}{240}$. Extracting gives $49+240=\boxed{289}.$

 


\end{enumerate}